\section{Анализ работ}

\subsection{Передача сообщений и многошаговые LoRa-сети}

Cardenas и соавторы~\cite{cardenas} описывают недорогую маломощную систему обмена сообщениями на LoRa. Эта работа близка по прикладной задаче: связь требуется пользователям, которым недоступна обычная инфраструктура. Для рассматриваемого исследования она задаёт сценарий коротких сообщений, однако не заменяет отдельной оценки стоимости маршрутизации.

Durand и Booysen~\cite{durand} исследуют mesh-топологию LoRa с помощью модели ns-3. Они анализируют доставку к шлюзу и энергопотребление при различных сценариях. Работа показывает необходимость совместного рассмотрения покрытия, ретрансляций и энергии. В отличие от схемы доставки к центральному шлюзу, здесь рассматривается обмен между парами равноправных узлов.

Udugampola и соавторы~\cite{udugampola} предлагают маршрутизацию с учётом положения и энергии узлов для подземных LoRa-сетей. В ней используются начальная фаза определения относительного положения, резервные ретрансляторы и переключение маршрутов. Результаты такой схемы зависят от геометрии и нагрузки подземной сети. Они мотивируют сравнение механизмов при различных условиях, но не дают универсального ранжирования протоколов для наземной оффлайн-связи.

Bor и соавторы~\cite{bor} экспериментально исследуют поведение LoRa и используют полученные зависимости в LoRaSim. Их подход важен для выбора модели потерь и обработки пересекающихся передач. Параметры распространения, измеренные в конкретной среде, при переносе в другую местность требуют калибровки.

\subsection{Маршрутизация и стоимость канала}

AODV обнаруживает маршрут по запросу и поддерживает его по мере необходимости~\cite{aodv}. OLSR заранее распространяет сведения о соседях и топологии, применяя многоточечные ретрансляторы MPR для сокращения служебного распространения~\cite{olsr}. Эти подходы позволяют сравнить затраты реактивного обнаружения и проактивного поддержания маршрутов.

Альтернативой служит управляемое лавинное распространение. Документация Meshtastic описывает задержку ретрансляции с учётом качества принятого сигнала и подавление части повторных передач после приёма дубликата~\cite{meshtastic}. В данной работе реализована упрощённая схема этого типа. Она воспроизводит выбранные механизмы распространения, но не весь стек и не все режимы прошивки Meshtastic.

Метрика ETX, предложенная De Couto и соавторами~\cite{etx}, учитывает ожидаемое число передач с повторениями и позволяет оценивать каналы иначе, чем по числу переходов. В классической постановке существенны вероятности доставки данных и подтверждения. В разработанном стенде используется упрощённая стоимость, пропорциональная времени передачи и обратно пропорциональная вероятности приёма кадра. Её ограничения при наличии подтверждений рассматриваются ниже.

Классические алгоритмы поиска кратчайших путей описаны, в частности, в~\cite{cormen}. Их сравнение характеризует вычисления на готовом графе. Для оценки LoRa mesh необходимо дополнительно учитывать, каким образом этот граф становится известен узлам и какой объём эфира расходуется на его поддержание.

Таким образом, в работе сопоставляются два вида затрат: обработка графа в памяти и получение сетевого состояния посредством радиопередач. Методика не предполагает, что самый быстрый алгоритм автоматически определяет лучший протокол, и проверяет эти уровни отдельно.
